\documentclass[exercice,a4]{thedocument}%
\usepackage{scratch3}%
\startdatakeys
    \source{25GENMATG11}
    \BO{2026}
    \cycle{4}
    \competencesDuSocle{CA,RA}
    \objectifsApprentissage{C4A1aa,C4A2ad,C4B2cb,C4D2aa,C4D3ac}
    \motsCles{probabilité,nombre premier,durée,Scratch,programme}
\enddatakeys
\begin{document}%
\medskip

Un digicode commande l'ouverture de la porte d'entrée de la maison de la grand-mère de Léna.

Léna a oublié le code. Elle sait qu'il est composé d'une lettre A, B, ou C, suivie d'un chiffre compris entre 0 et 9.

\medskip

\begin{enumerate}
\item Proposer deux codes différents que Léna peut tester.
\item Quelle est la probabilité que la grand-mère de Léna ait choisi la lettre C dans son code ?
\item Montrer que la probabilité que la grand-mère de Léna ait choisi le chiffre 7 dans son code est $\dfrac{1}{10}$.
\item Léna se souvient que sa grand-mère, enseignante de mathématiques à la retraite, aime bien les nombres premiers. Quelle est la probabilité que le code choisi par sa grand-mère comporte un nombre premier ?
\item 
	\begin{enumerate}
		\item Léna décide de tester tous les codes possibles. Elle estime qu'il lui faut 5 secondes pour essayer un code. Réussira-t-elle à ouvrir la porte de la maison en moins de 3 minutes?
		\item Le format de ce code garantit-il la sécurité de la maison? Comment pourrait-on améliorer ce système de code?
	\end{enumerate}
\item Chaque fois qu'un utilisateur saisit un code, un programme lui annonce si le code est correct ou faux. Le programme utilisé est noté ci-dessous.

\begin{center}
\begin{scratch}
\blockinit{quand \greenflag est cliqué}
\blocksensing{demander \ovalnum{Saisir une lettre parmi A, B et C} et attendre}
\blockvariable{mettre \selectmenu{lettre saisie} à \ovaloperator{\ovalvariable{réponse}}}
\blocksensing{demander \ovalnum{Saisir un nombre entre 0 et 9} et attendre}
\blockvariable{mettre \selectmenu{chiffre saisi} à \ovaloperator{\ovalvariable{réponse}}}
\blockifelse{si \booloperator{\ovalmove{lettre saisie} = \ovalnum{B} et \ovalmove{chiffre saisi} = \ovalnum{7}} alors}
{\blocklook{dire \ovalnum{Code correct} pendant \ovalnum{2} secondes}}
{\blocklook{dire \ovalnum{Code faux} pendant \ovalnum{2} secondes}}
\end{scratch}
\end{center}

	\begin{enumerate}
		\item Léna saisit le code B5. Qu'affiche le programme ?
		\item D'après ce programme, quel est le code qui permet d'entrer dans l'immeuble de la grand-mère
de Léna ?
	\end{enumerate}
\end{enumerate}
\showthesource
\end{document}
